%% C4 nivel 2 — Diagrama de contenedores (plantilla)
%% Compilar: tectonic -X compile c4-contenedores.tex
\documentclass[11pt,a4paper]{article}
\usepackage{../../docs/latex/legasoft}
\usepackage{../../docs/latex/legasoft-diagramas}

\lgtitulo{C4 nivel 2: diagrama de contenedores}
\lgtituloCorto{C4 — Contenedores}
\lgcodigo{LGS-MOD-C4-2}
\lgversion{0.1}
\lgestado{Plantilla}
\lgfecha{\campo{fecha}}
\lgautor{\lgArquitecto\ (\lgArquitectoRol)}

\begin{document}

\begin{center}
{\sffamily\bfseries\LARGE\color{lgPrimario} C4 nivel 2: diagrama de contenedores\par}
\vspace{2pt}
{\sffamily\normalsize\color{lgGris} \campo{nombre del sistema} · Legasoft\par}
\end{center}
\vspace{6pt}
{\color{lgBorde}\hrule height 0.8pt}
\vspace{10pt}

\begin{porcompletar}[Cómo usar esta plantilla]
El nivel 2 abre el sistema y muestra sus unidades desplegables —aplicaciones,
servicios y almacenes de datos— con la tecnología de cada una. Las tecnologías
todavía no están decididas: mientras lo estén, deje los campos en ámbar y
registre cada elección como ADR en \ruta{docs/arquitectura/adr/}.
\end{porcompletar}

% =============================================================================
\section*{Diagrama}

\begin{figure}[H]
\centering
\ajustaancho{%
\begin{tikzpicture}[c4]

  % --- Personas -------------------------------------------------------------
  \node[c4persona, text width=3.0cm, align=center, minimum height=1.5cm] (usuario)
    {\ctexto{\campoclaro{Usuario}}{}{Usa la aplicación.}};

  % --- Contenedores dentro del límite del sistema ---------------------------
  \node[c4contenedor, below=2.0cm of usuario] (web)
    {\ctexto{Aplicación web}{\campo{tecnología}}{Presenta la interfaz y consume la API.}};

  \node[c4contenedor, below=1.7cm of web] (api)
    {\ctexto{API de aplicación}{\campo{tecnología}}{Expone los contratos y ejecuta la lógica de negocio.}};

  \node[c4datos, right=2.4cm of api] (bd)
    {\ctexto{Base de datos}{\campo{gestor}}{Persiste la información del dominio.}};

  % Límite del sistema
  \begin{scope}[on background layer]
    \node[c4limite, fit=(web)(api)(bd),
          label={[font=\sffamily\scriptsize\itshape, text=lgPrimario,
                  anchor=north west]north west:{\campo{Nombre del sistema}}}]
      (limite) {};
  \end{scope}

  % --- Sistema externo ------------------------------------------------------
  \node[c4externo, left=2.4cm of api] (ext)
    {\ctexto{\campo{Sistema externo}}{}{\campo{Qué provee.}}};

  % --- Relaciones -----------------------------------------------------------
  \draw[c4rel] (usuario.south) -- (web.north)
    node[c4etq, pos=0.5] {\campo{opera}\\{\scriptsize[HTTPS]}};

  \draw[c4rel] (web.south) -- (api.north)
    node[c4etq, pos=0.5] {llama\\{\scriptsize[JSON/HTTPS]}};

  \draw[c4rel] (api.east) -- (bd.west)
    node[c4etq, pos=0.5] {lee y escribe\\{\scriptsize\campo{protocolo}}};

  \draw[c4relbi] (api.west) -- (ext.east)
    node[c4etq, pos=0.5] {integra\\{\scriptsize\campo{API}}};

\end{tikzpicture}}
\caption{Contenedores del sistema, su tecnología y sus relaciones.}
\end{figure}

\cleyenda

% =============================================================================
\Needspace*{20\baselineskip}
\section*{Contenedores}

\begin{center}
\begin{tabularx}{\linewidth}{@{}L{3.4cm} L{2.8cm} Y@{}}
\toprule
\sffamily\bfseries\color{lgPrimario}Contenedor &
\sffamily\bfseries\color{lgPrimario}Tecnología &
\sffamily\bfseries\color{lgPrimario}Responsabilidad \\
\midrule
Aplicación web     & \campo{por definir} & Interfaz de usuario; no contiene reglas de negocio. \\
API de aplicación  & \campo{por definir} & Contratos, validación y lógica de negocio. \\
Base de datos      & \campo{por definir} & Persistencia e integridad de los datos del dominio. \\
\bottomrule
\end{tabularx}
\captionof{table}{Contenedores y su responsabilidad.}
\end{center}

\begin{nota}[Contratos entre contenedores]
La definición detallada de los contratos entre la aplicación web y la API se
documenta en \texttt{LGS-ARQ-001}. Esos contratos son la base sobre la que el
\lgQARol\ construye las pruebas de integración de \ruta{tests/integration}.
\end{nota}

\end{document}
